\section*{Chapter \ref{ch:b1:fractions} --- Rational Fractions}

\begin{solution}{exo:b1:fractions:1}
$\dfrac{1}{X^2 - 1}$: simple poles $\pm 1$; cover-up:
$\dfrac{1/2}{X-1} - \dfrac{1/2}{X+1}$.

$\dfrac{X}{X^2 - 3X + 2} = \dfrac{X}{(X-1)(X-2)}$: cover-up gives
$\frac{1}{1-2} = -1$ at $1$ and $\frac{2}{2-1} = 2$ at $2$:
$\dfrac{-1}{X-1} + \dfrac{2}{X-2}$.

$\dfrac{X^2+1}{X(X-1)}$: the degree is $0$, so there is an integer
part: dividing, $X^2 + 1 = (X^2 - X) + (X + 1)$, so $F = 1 +
\frac{X+1}{X(X-1)}$. Cover-up on the remainder: $\frac{1}{-1} = -1$
at $0$, $\frac{2}{1} = 2$ at $1$:
\[
F = 1 - \frac{1}{X} + \frac{2}{X-1} .
\]
\end{solution}

\begin{solution}{exo:b1:fractions:2}
$\dfrac{1}{X(X^2+1)}$: shape $\frac aX + \frac{bX + c}{X^2 + 1}$.
Cover-up at $0$: $a = 1$. Limit of $xF$: $0 = a + b$, so $b = -1$.
Evaluation at $X = 1$: $\frac12 = 1 + \frac{c - 1}{2}$, so $c = 0$:
\[
\frac{1}{X(X^2+1)} = \frac 1X - \frac{X}{X^2+1} .
\]

$\dfrac{X^3}{X^2+X+1}$: division: $X^3 = (X^2+X+1)(X - 1) + 1$, so
\[
\frac{X^3}{X^2+X+1} = X - 1 + \frac{1}{X^2 + X + 1} ,
\]
already in real decomposed form (the quadratic has negative
discriminant).
\end{solution}

\begin{solution}{exo:b1:fractions:3}
$\dfrac{1}{X^2(X-1)}$: shape $\frac{a}{X^2} + \frac bX +
\frac{c}{X-1}$. Cover-up at the double pole $0$:
$a = \bigl[\frac{1}{X-1}\bigr]_{0} = -1$. Cover-up at $1$: $c = 1$.
Limit of $xF$: $0 = b + c$, so $b = -1$:
\[
\frac{1}{X^2(X-1)} = -\frac{1}{X^2} - \frac1X + \frac{1}{X-1} .
\]

$\dfrac{X+1}{(X-1)^3}$: with $Y = X - 1$, the numerator is $Y + 2$:
\[
\frac{Y + 2}{Y^3} = \frac{1}{Y^2} + \frac{2}{Y^3}
= \frac{1}{(X-1)^2} + \frac{2}{(X-1)^3} .
\]
\end{solution}

\begin{solution}{exo:b1:fractions:4}
The poles are the fourth roots of unity $1, \iu, -1, -\iu$, all
simple. Simple-pole formula with $B' = 4X^3$: coefficient at $a$ is
$\frac{1}{4a^3} = \frac{a}{4a^4} = \frac a4$ (using $a^4 = 1$). So,
over $\C$:
\[
\frac{1}{X^4 - 1}
= \frac{1/4}{X - 1} - \frac{1/4}{X + 1}
+ \frac{\iu/4}{X - \iu} - \frac{\iu/4}{X + \iu} .
\]
Grouping the conjugate pair (common denominator $X^2 + 1$):
$\frac{\iu}{4}\bigl(\frac{1}{X-\iu} - \frac{1}{X+\iu}\bigr) =
\frac{\iu}{4}\cdot\frac{2\iu}{X^2+1} = \frac{-1/2}{X^2+1}$. Over
$\R$:
\[
\frac{1}{X^4 - 1}
= \frac{1/4}{X-1} - \frac{1/4}{X+1} - \frac{1/2}{X^2 + 1} .
\]
\emph{Check at $X = 0$:} $-1 = -\frac14 - \frac14 - \frac12$.
\end{solution}

\begin{solution}{exo:b1:fractions:5}
$\dfrac{X^2}{(X^2+1)^2} = \dfrac{(X^2 + 1) - 1}{(X^2+1)^2} =
\dfrac{1}{X^2+1} - \dfrac{1}{(X^2+1)^2}$.

Hence a primitive:
\[
\int \frac{x^2\,\dd x}{(x^2+1)^2}
= \arctan x - \frac12\Bigl(\arctan x + \frac{x}{x^2+1}\Bigr) + C
= \frac12\arctan x - \frac{x}{2(x^2+1)} + C .
\]
\end{solution}

\begin{solution}{exo:b1:fractions:6}
The poles $0, -1, \dots, -n$ are simple. Cover-up at $-k$:
\[
c_k = \frac{n!}{\prod_{j \neq k} (-k + j)}
= \frac{n!}{\bigl(\prod_{j=0}^{k-1}(j - k)\bigr)
\bigl(\prod_{j=k+1}^{n}(j-k)\bigr)}
= \frac{n!}{(-1)^k k!\,(n-k)!} = (-1)^k \binom nk .
\]
So
\[
\frac{n!}{X(X+1)\cdots(X+n)}
= \sum_{k=0}^{n} \frac{(-1)^k \binom nk}{X + k} .
\]
(Sanity check for $n = 1$: $\frac{1}{X(X+1)} = \frac1X -
\frac{1}{X+1}$.)
\end{solution}

\begin{solution}{exo:b1:fractions:7}
$P = X^n - 1$ has the $n$ simple roots $\omega^k$
(\cref{thm:b1:complex:roots}), so the logarithmic-derivative identity
of \cref{exo:b1:poly:11} reads
\[
\sum_{k=0}^{n-1} \frac{1}{X - \omega^k}
= \frac{P'(X)}{P(X)} = \frac{n X^{n-1}}{X^n - 1} .
\]
At $X = 2$, $n = 4$: right side $= \frac{4 \times 8}{15} =
\frac{32}{15}$. Left side: $\frac{1}{2-1} + \frac{1}{2+1} +
\frac{1}{2 - \iu} + \frac{1}{2 + \iu} = 1 + \frac13 + \frac{4}{5} =
\frac{15 + 5 + 12}{15} = \frac{32}{15}$, using $\frac{1}{2-\iu} +
\frac{1}{2+\iu} = \frac{4}{5}$.
\end{solution}

\begin{solution}{exo:b1:fractions:8}
Cover-up: $\dfrac{1}{k(k+1)(k+2)} = \dfrac{1/2}{k} - \dfrac{1}{k+1} +
\dfrac{1/2}{k+2}$. Rewrite as a telescoping difference:
\[
\frac{1}{k(k+1)(k+2)}
= \frac12\Bigl(\frac{1}{k(k+1)} - \frac{1}{(k+1)(k+2)}\Bigr),
\]
(expand to check --- or subtract the two decompositions). Summing:
\[
S_n = \frac12\Bigl(\frac{1}{1 \times 2} - \frac{1}{(n+1)(n+2)}\Bigr)
= \frac14 - \frac{1}{2(n+1)(n+2)}
\xrightarrow[n \to \infty]{} \frac14 .
\]
\end{solution}

\begin{solution}{exo:b1:fractions:9}
Decompose, for $0 \leq m \leq n - 1$, the fraction $\frac{X^m}{P}$
(degree $m - n \leq -1$, poles simple): the simple-pole formula gives
\[
\frac{X^m}{P} = \sum_{i=1}^{n} \frac{x_i^m / P'(x_i)}{X - x_i} .
\]
Multiply by $X$ and let $X \to +\infty$: the left side tends to the
limit of $X^{m+1}/P$, which is $0$ if $m \leq n - 2$ and $1$ if $m =
n - 1$ ($P$ is monic of degree $n$); the right side tends to
$\sum_i \frac{x_i^m}{P'(x_i)}$. Hence
\[
\sum_{i=1}^{n} \frac{x_i^{m}}{P'(x_i)} =
\begin{cases}
0 & \text{for } 0 \leq m \leq n-2,\\
1 & \text{for } m = n - 1,
\end{cases}
\]
which contains both announced identities ($m = 0$ requires $n \geq
2$). (Interpretation via \cref{thm:b1:poly:lagrange}: these sums are
the leading coefficients of the Lagrange interpolants of $X^m$, and
interpolating a polynomial of degree $\leq n-1$ at $n$ points
reproduces it exactly.)
\end{solution}
